\chapter{Retain state between lossless video frames}
\MARSChapterHeading{Retain state between lossless video frames}
\label{revision:lossless-video}

A sequence of PNG images can represent a video, but it describes each image
independently. Often only a small part of the next picture changes. We can
store those changes and retain the preceding decoded image as state. The
saving comes with a new responsibility: a delta has meaning only in relation
to the correct previous frame.

This chapter develops MLV1, the book's small lossless teaching format. Every
picture is 64 by 64 pixels. Some records contain complete RGB images; others
contain XOR changes. A directory supports seeking from an independent image,
and a player gives accepted pictures their display intervals. The complete
\texttt{lossless-video.asm} in \texttt{revision/labs/files} contains the codec
and player. The original construction package remains available for building
all its entries; the independent lab locates the rectangle containing changes.

This work assumes the earlier native RGB words, CRC, bounded buffers, file I/O
and MSA chapters. MLV1 is a local format, not a stream that a general-purpose
video player is expected to recognize. The following H.263 chapters will apply
related state reasoning to a restricted interoperable codec.

\section{Give the previous frame an exact meaning}

One serialized image contains $4096\cdot3=12288$ RGB bytes. Its working image
contains 4096 native \texttt{0x00RRGGBB} words, or 16384 bytes, with a packed
64-word row stride. Load a native word numerically and extract its channels
when writing RGB. The file bytes for word \texttt{0x00123456} are
\texttt{12 34 56} in either target byte order.

For corresponding previous and current words $P_i,C_i$, form
\[
 D_i=P_i\mathbin{\oplus}C_i.
\]
Then $P_i\mathbin{\oplus}D_i=C_i$. XOR reverses itself. A zero delta preserves
the word; a nonzero delta marks a changed pixel. After accepting the transition,
the next delta must use $C$, not the image that preceded $P$. Reusing an older
full image would lose every intervening change that the new delta leaves zero.

The scalar and MSA XOR entries receive two word arrays, an output and a count
zero through 4096. Output may be disjoint or exactly equal to an input. Load
both corresponding inputs before storing, making those exact aliases safe.
Partial overlap is outside this contract. Count zero succeeds without any
pointer access. For positive counts, word alignment suffices; a vector route
must not silently require sixteen-byte alignment.

A complete vector group performs four XOR words:
\Needspace{9\baselineskip}
\MARSCodeLabel{Compute four independent pixel differences}
\begin{minted}{mars-idea}
    ld.w $w0, 0($a0)
    ld.w $w1, 0($a1)
    xor.v $w2, $w0, $w1
    st.w $w2, 0($a2)
\end{minted}
The function also returns the number of nonzero words. Compare all 32 bits
with zero; hexadecimal \texttt{80000000} is nonzero even though a signed word
interpretation is negative. This numerical kernel accepts all word patterns,
whereas serialized image colors require a zero high byte. When it compares
previous and current images, its count is changed pixels. When it applies XOR
data to an image, the same count describes nonzero results instead. Name the
meaning at the call site.

\section{Describe a frame record without losing its boundaries}

The 24-byte file header begins with ASCII \texttt{MLV1}, then big-endian
unsigned width and height of two bytes each. Both must be 64. A four-byte frame
count admits one through 256; a two-byte full-frame period admits one through
64; two flag bytes are zero. A four-byte duration admits one through 1000
milliseconds, and the final four reserved bytes are zero. These bounds must be
established before indexing, allocating or dividing by the period.

Each record has a sixteen-byte header: kind at byte zero, three zero reserved
bytes, four-byte sequence at offset four, payload count at eight, and CRC-32
at twelve. Kind zero means a full RGB image and kind one a delta. Sequence
starts at zero and advances exactly once per record. The CRC covers the first
twelve header bytes followed by the payload, excluding its own four bytes.
There is no implicit padding between records.

A full payload is exactly 12288 bytes. A delta consists of commands, each with
a big-endian two-byte skip and two-byte changed count followed by three XOR
bytes per changed pixel. Begin the pixel cursor at zero for every delta. Skip
leaves words unchanged; changed triples XOR into the following words. Each
command must advance by at least one pixel, and both pixel and payload extents
must finish exactly. A zero/zero command is invalid. A final nonzero skip with
no changed pixels is valid and needs no terminator.

For the one-pixel update, the payload is
\texttt{00 00 00 01 01 00 00 0f ff 00 00}. Its first command skips zero and
changes one pixel using red XOR one. The cursor is now one. The second command
skips 4095 and changes zero, finishing at 4096. Eleven payload bytes describe
this transition. The three bytes \texttt{01 00 00} are a difference, not the
new pixel's complete color.

Each command has four header bytes, and there are at most 4096 commands and
4096 changed triples. Thus the conservative delta payload bound is
$4\cdot4096+3\cdot4096=28672$ bytes. The complete-file bound is
$24+256(16+28672)=7344152$ bytes. These bounds admit even inefficient valid
command streams; they do not prescribe an encoder's output size.

\section{Validate the update before publishing it}

\texttt{video\_apply\_record} takes one exact record window, a working image
and a native context. The context supplies image capacity, expected sequence
and period. Successful publication advances the expected sequence and records
the accepted kind and 16384 active image bytes. Failure returns one and zero,
retaining every image and context byte. Input bytes remain immutable.

First admit alignment, capacities, header extent, kind, reserved fields,
sequence, required full-frame positions, payload length and CRC. Then walk a
delta without writing pixels. After each skip, compare the changed count with
the remaining pixels, and compare its required triples with remaining payload
bytes. If all checks succeed, replay the immutable commands to apply them.
The second pass uses boundaries already established by the first.

Why pay for two walks? Suppose the first command changes a visible pixel but
the final command is zero/zero. A decoder that applies while validating has
already damaged the displayed image when it discovers the final error. A valid
CRC does not make those commands meaningful. The accepted and rejected native
observations below use the same initial image but different record validity.
\begin{figure}[!ht]
\centering
\begin{tikzpicture}[alt={A native Memory row starts with four a5 guard bytes. An outline marks c2 61 01 00, the first active word 000161c2 after a red XOR of one.}]
\node[anchor=north west,inner sep=0,font=\small\bfseries] (title) at (0,0) {A validated record changes the first pixel};
\node[anchor=north west,inner sep=0] (a) at ([yshift=-2mm]title.south west) {\MARSRevisionCrop{revision/captures/lossless/lossless-record.png}{0.99\linewidth}{0.03125}{0.33888889}{0.58203125}{0.60888889}};
\begin{scope}[shift={(a.south west)},x={($(a.south east)-(a.south west)$)},y={($(a.north west)-(a.south west)$)}]
\draw[black,line width=.7pt] (0.391250,0.005000) rectangle (0.582500,0.443809);
\end{scope}
\end{tikzpicture}
\caption{The outlined first active word becomes 000161c2 from 000061c2 after red XOR one. The preceding four a5 guard bytes remain unchanged; the following words belong to unchanged pixels.}
\label{fig:revision-lossless-native-record}
\end{figure}

\begin{figure}[!ht]
  \centering
  \begin{tikzpicture}[alt={The same native Memory interval after a checksummed record ends in a zero-zero command. The outlined first active word remains c2 61 00 00, or 000061c2.}]
    \node[anchor=north west,inner sep=0,font=\small\bfseries] (title) at (0,0) {The late invalid command publishes no pixel};
    \node[anchor=north west,inner sep=0] (a) at ([yshift=-2mm]title.south west) {\MARSRevisionCrop{revision/captures/lossless/lossless-rejected-record.png}{0.99\linewidth}{0.03125}{0.33888889}{0.58203125}{0.60888889}};
    \begin{scope}[shift={(a.south west)},x={($(a.south east)-(a.south west)$)},y={($(a.north west)-(a.south west)$)}]
      \draw[black,line width=.7pt] (0.391250,0.005000) rectangle (0.582500,0.443809);
    \end{scope}
  \end{tikzpicture}
  \caption{The outlined word remains 000061c2 when the final command makes no progress. Its record CRC was repaired before the run, so rejection tests command validity rather than a checksum mismatch.}
  \label{fig:revision-lossless-native-rejected-record}
\end{figure}


At sequence indices divisible by the positive period, require a full record.
Other indices may also use full records. A full frame replaces previous pixels
and establishes a fresh prediction state. A delta requires the immediately
preceding accepted image. Preserve saved registers, frame pointer and entry
stack pointer throughout; nested codec calls use a sixteen-byte-aligned stack.
Temporary frames and contexts belong to the current call, not a shared cursor
that another playback session might change.

\section{Choose an encoding after measuring its byte cost}

Start with a full encoder. Validate all 4096 source words, including the last
high byte, and admit output capacity before writing. Emit RGB channels and a
complete header/CRC, then update the encoder's previous image and next sequence.
If validation or capacity fails, record output, previous image and context
remain unchanged. An uncommitted source image must not become the prediction
state for a later retry.

For a nonperiodic frame, compute XOR words in private scratch. Walk maximal
unchanged runs and changed runs. If $k$ commands cover $m$ changed pixels,
the planned payload is $4k+3m$ bytes. The changed count alone omits the command
cost. An alternating image can have fewer changed pixels than a uniformly
changed image yet require more command headers.

\begin{table}[!ht]
\centering
\begin{tabular}{l r r r l}
\toprule
Transition & Changes & Commands & Delta payload & Worked choice \\
\midrule
Unchanged & 0 & 1 & 4 & Delta \\
Only final pixel & 1 & 1 & 7 & Delta \\
Only first pixel & 1 & 2 & 11 & Delta \\
Two-by-two patch & 4 & 3 & 24 & Delta \\
Every pixel & 4096 & 1 & 12292 & Full \\
\bottomrule
\end{tabular}
\caption{Payload cost includes command headers, in addition to changed triples.}
\end{table}

The worked encoder chooses full when the planned delta is at least 12288 bytes,
including a tie. Otherwise it checks room for sixteen header bytes plus that
payload, then repeats the run walk to emit it. A mandatory full position takes
precedence over this comparison. Equivalent encoders may choose additional
full frames or different valid command boundaries. A changed span may even
contain a zero XOR triple. Correct decoding and format constraints determine
validity; a particular greedy partition is not the file format.

The complete-file encoder owns different failure boundaries from the single
record encoder. It retains immutable packed source images and publishes its
file count and record count only on complete success. Its staging bytes,
previous-image scratch and XOR scratch may already have changed when a later
frame fails. Keep any old animation in separate storage if it must survive an
attempted replacement.

\section{Index the whole file, then seek from a full frame}

The moving fixture has period four and eight frames. Its record directory is:
\Needspace{16\baselineskip}
\begin{table}[!ht]
\centering
\begin{tabular}{r l r r}
\toprule
Sequence & Kind & File offset & Payload bytes \\
\midrule
0 & Full & 24 & 12288 \\
1 & Delta & 12328 & 11 \\
2 & Delta & 12355 & 24 \\
3 & Delta & 12395 & 4 \\
4 & Full & 12415 & 12288 \\
5 & Delta & 24719 & 17 \\
6 & Delta & 24752 & 7 \\
7 & Delta & 24775 & 4 \\
\bottomrule
\end{tabular}
\caption{The final record ends at file byte count 24795. Directory entries are offsets.}
\end{table}

\texttt{video\_index} validates the global header and every record, including
sequence, periodic full positions and exact file completion. Its scratch image
can change during validation. The public directory and context remain unchanged
on failure. A second pass can publish offsets only after the whole immutable
file is accepted. One offset occupies eight native bytes; it becomes an address
only after adding the retained file base.

A seek receives a validated file and directory and a target index. The displayed
image could belong to an unrelated position, so reconstruct in separate scratch.
Scan backward from the target to the latest full record, apply it and every
following delta through the target, then copy the complete image to display
and publish target, replay start and replay count. Failure retains the old
destination and context, though scratch may have changed.

For target six, the moving fixture replays four, five and six. An equivalent
fixture that stores frame five as an extra full image replays only five and
six, yielding identical pixels. Computing only the required periodic boundary
would miss that nearer full frame. Period bounds how far backward a valid seek
must look; it does not identify every full frame that an encoder may choose.
\begin{figure}[!ht]
\centering
\begin{tikzpicture}[every node/.style={font=\small,align=center},alt={The ordinary moving file reconstructs target six using full frame four, delta five and delta six. An equivalent file with an additional full frame five reconstructs the same target using only full five and delta six.}]
\node[draw,minimum width=20mm,minimum height=10mm] (a) at (0,0) {4\\Full};
\node[draw,minimum width=20mm,minimum height=10mm] (b) at (3,0) {5\\Delta};
\node[draw,minimum width=20mm,minimum height=10mm] (c) at (6,0) {6\\Delta};
\draw[-{Latex}] (a)--(b);\draw[-{Latex}] (b)--(c);
\node[above=2mm of b] {Three records for the moving file};
\node[draw,minimum width=20mm,minimum height=10mm] (d) at (3,-1.6) {5\\Full};
\node[draw,minimum width=20mm,minimum height=10mm] (e) at (6,-1.6) {6\\Delta};
\draw[-{Latex}] (d)--(e);
\node[align=left] at (.1,-1.6) {An additional\\full frame};
\end{tikzpicture}
\caption{Two valid encodings reach identical target pixels through different replay starts.}
\label{fig:revision-lossless-replay}
\end{figure}

\begin{figure}[!ht]
  \centering
  \begin{tikzpicture}[alt={Native seek-context results show target six, replay start four and three applied records. The replay-start field has a dashed outline.}]
    \node[anchor=north west,inner sep=0,font=\small\bfseries] (title) at (0,0) {Target six starts at full frame four};
    \node[anchor=north west,inner sep=0] (a) at ([yshift=-2mm]title.south west) {\MARSRevisionCrop{revision/captures/lossless/lossless-seek.png}{0.99\linewidth}{0.03125}{0.31555556}{0.58203125}{0.60888889}};
    \begin{scope}[shift={(a.south west)},x={($(a.south east)-(a.south west)$)},y={($(a.north west)-(a.south west)$)}]
      \draw[black,line width=.7pt] (0.582500,0.313824) rectangle (0.965000,0.614647);
      \draw[black,dashed,line width=.7pt] (0.200000,0.005000) rectangle (0.582500,0.305824);
      \draw[black,line width=.7pt] (0.582500,0.005000) rectangle (0.965000,0.305824);
    \end{scope}
  \end{tikzpicture}
  \caption{The three outlined native fields publish target six, replay start four (dashed), and three replayed records. The result follows full frame four and deltas five and six; it does not depend on the earlier display image.}
  \label{fig:revision-lossless-native-seek}
\end{figure}


Loading and indexing are separate promises. A loader establishes a complete
nonempty byte count using actual read results and EOF. Indexing establishes
that those bytes form this format. After storing and reloading a file, validate
the reloaded bytes afresh before trusting any earlier directory. Likewise,
overwriting staging storage invalidates an old directory referring to it.
\texttt{video\_store} is a transport routine, not an atomic file replacement:
a write failure can leave a truncated output. Validate staged content before
opening its output path.

\section{Keep logical display intervals separate from execution time}

An unchanged frame still occupies an interval. For the moving file, eight
frames at forty milliseconds specify 320 milliseconds of logical duration.
That does not promise the host or IDE will finish or repaint in 320 milliseconds.
The player validates input before reading its initial clock anchor or changing
display pixels, then seeks independently to its start frame.

Let the initial anchor be 1000 and duration forty. The first deadline is 1040,
the second 1080, and the eighth 1320. Advance this original schedule after each
accepted frame, even when an earlier one was late. If time is below the current
deadline, request only the positive remaining interval and read time again.
An early wakeup needs another check. Equal time is on time; later time records
lateness once for that interval. This player does not drop frames or re-anchor.

For readings 1010, 1035 and 1040 after the anchor, sleep requests are thirty
and five milliseconds. For readings 1030 then 1047, request ten and record
seven milliseconds late. A later reading below the preceding observed time
is an error. Deadlines and readings are full unsigned doublewords; reject
an overflowing deadline before presenting that interval.

Service thirty returns time halves in two registers. The helper zero-extends
the low half before joining it to the shifted high half:
\Needspace{8\baselineskip}
\MARSCodeLabel{Join the two time-service result halves}
\begin{minted}{mars-idea}
    dsll $t0, $a0, 32
    dsrl $t0, $t0, 32
    dsll $v1, $a1, 32
    or $v1, $v1, $t0
\end{minted}
A low half with bit 31 set must not fill the result's upper half with sign bits.
The scripted grading clock checks the protocol without host waiting: one anchor,
a reading after each accepted frame and a reading after each sleep request.
Actual Debug observations use the host clock. The final native context below
separates accepted count and logical duration from observed lateness.
\begin{figure}[!ht]
  \centering
  \begin{tikzpicture}[alt={Native playback fields show eight accepted frames, 8 late intervals, maximum lateness 104 milliseconds and logical duration 320 milliseconds. Lateness is from this actual Debug run.}]
    \node[anchor=north west,inner sep=0,font=\small\bfseries] (title) at (0,0) {Logical duration and host lateness have separate fields};
    \node[anchor=north west,inner sep=0] (a) at ([yshift=-2mm]title.south west) {\MARSRevisionCrop{revision/captures/lossless/lossless-playback.png}{0.99\linewidth}{0.03125}{0.31555556}{0.58203125}{0.60888889}};
    \begin{scope}[shift={(a.south west)},x={($(a.south east)-(a.south west)$)},y={($(a.north west)-(a.south west)$)}]
      \draw[black,line width=.7pt] (0.200000,0.313824) rectangle (0.582500,0.614647);
      \draw[black,dashed,line width=.7pt] (0.582500,0.313824) rectangle (0.965000,0.614647);
      \draw[black,line width=.7pt] (0.200000,0.005000) rectangle (0.582500,0.305824);
      \draw[black,line width=.7pt] (0.582500,0.005000) rectangle (0.965000,0.305824);
    \end{scope}
  \end{tikzpicture}
  \caption{This captured Debug run accepts eight frames in 320 milliseconds of logical intervals. It records 8 late intervals (dashed field) and maximum lateness 104 milliseconds. These are host-clock observations for this run, not a guarantee about later runs or screen presentation time.}
  \label{fig:revision-lossless-native-playback}
\end{figure}


Only complete playback success publishes the four final timing/result fields.
A failure returns the actual number of already accepted frames and leaves
those images visible; it does not roll them back. Its final context fields
remain unchanged. This differs from a failed seek, which publishes no new
image. State the guarantee for each operation rather than assuming every
function has one universal failure transaction.

\section{Run in a fresh directory and inspect a transition}

The worked caller encodes its native source frames, indexes the staged file,
checks an independent seek, writes \texttt{animation.mlv}, reloads it and plays
the validated result. Run in a fresh working directory: the program opens that
output for writing. Here \texttt{CLI} names the full path of the built MARS JAR.
\Needspace{8\baselineskip}
\MARSCodeLabel{Run the complete file codec and player}
\begin{Verbatim}
java -jar "$CLI" --start-at-main \
  /path/to/mars-book/revision/labs/files/lossless-video.asm
\end{Verbatim}
Success prints \texttt{Video file loaded and played.} In IntelliJ, select the
caller for MARS Run or Debug and use the same fresh working directory. Connect
a 64-column bitmap to the actual \texttt{playback\_frame+4} address when
following presentation. A completed store establishes simulator state; it does
not establish when the screen repainted.

The fixture writer can also create an independently planned moving file and
PPM previews in a new output folder:
\Needspace{8\baselineskip}
\MARSCodeLabel{Create an independently planned animation for inspection}
\begin{Verbatim}
python3 packages/034-lossless-video/construction/video-fixtures/fixture.py \
  moving --output moving-input
\end{Verbatim}
That file can have a different valid representation from an equivalent encoder's
output. Compare reconstructed pixels and valid record structure rather than
requiring identical bytes from every solution. The lab tests an additional full
record, malformed late commands and retries before assigning changed bounds.

\Needspace{30\baselineskip}
\section{Locate changes without confusing count and area}

Implement \texttt{video\_change\_bounds} in the independent file. It compares
two native word images and returns the smallest axis-aligned rectangle
containing their changed pixels, with exclusive right and bottom coordinates,
plus the changed-word count. An unchanged pair returns four zero bounds and
count zero. This describes a region; it does not encode a delta or replace the
format's sequence and validation rules.

The moving two-by-two patch occupies $x=5,6$ and $y=5,6$, giving bounds
$(5,5,7,7)$ and count four. If only opposite corners of a 64-square picture
change, the rectangle covers all 4096 positions while the changed count is two.
A rectangular repaint region, a run-coded payload and a changed count answer
different questions about the same transition. The independent interface also
accepts smaller dimensions for checking. Its two-by-three demo changes two
separated words whose containing rectangle covers six cells.
\begin{figure}[!ht]
\centering
\begin{tikzpicture}[every node/.style={font=\small},alt={A two-column, three-row image has changes at column one row zero and column zero row two. Their smallest containing rectangle spans all six cells. Its exclusive bounds are zero, zero, two, three, but only two words changed.}]
\fill[black!12] (1.2,0) rectangle (2.4,-1.2);
\fill[black!12] (0,-2.4) rectangle (1.2,-3.6);
\draw[step=1.2] (0,0) grid (2.4,-3.6);
\draw[dashed,line width=.9pt] (-.05,.05) rectangle (2.45,-3.65);
\node at (1.8,-.6) {$\times$};\node at (.6,-3) {$\times$};
\node at (.6,.3) {0};\node at (1.8,.3) {1};\node at (1.2,.8) {column};
\node at (-.4,-.6) {0};\node at (-.4,-1.8) {1};\node at (-.4,-3) {2};
\node[align=left] at (5,-1.8) {Changed words: 2\\Rectangle area: 6\\Bounds: $(0,0,2,3)$};
\end{tikzpicture}
\caption{Exclusive ends lie beyond the last selected column and row. The containing rectangle also includes unchanged cells.}
\label{fig:revision-lossless-change-bounds}
\end{figure}


Choose a single scan with explicit coordinates, or separate passes that first
locate changed rows and then their horizontal extent. Both must compare original
corresponding words and preserve inputs. The lab specifies dimensions, native
alignment, exact result publication, aliases and invalid-input behavior. The
full decoder and encoder construction remains a separate capstone route.
